\documentclass[11pt,a4paper]{article}

% ============================================================
% Basic typography and page layout
% ============================================================
\usepackage[T1]{fontenc}
\usepackage[utf8]{inputenc}
\usepackage{lmodern}
\usepackage[a4paper,margin=2.5cm]{geometry}
\usepackage{microtype}
\usepackage{setspace}
\setstretch{1.02}

% ============================================================
% Mathematics and physics
% ============================================================
\usepackage{amsmath,amssymb,amsfonts,mathtools}
\usepackage{bm}
\usepackage{physics}
\usepackage{siunitx}
\sisetup{separate-uncertainty=true,per-mode=symbol}

% ============================================================
% Figures, tables and captions
% ============================================================
\usepackage{graphicx}
\usepackage{subcaption}
\usepackage{booktabs}
\usepackage{multirow}
\usepackage{array}
\usepackage{longtable}
\usepackage{float}
\usepackage{placeins}
\usepackage[font=small,labelfont=bf]{caption}
\usepackage{tablefootnote}

% ============================================================
% Lists, code and colors
% ============================================================
\usepackage{enumitem}
\usepackage{xcolor}
\usepackage[table]{xcolor}
\usepackage{transparent}
\usepackage{listings}
\lstset{
  basicstyle=\ttfamily\small,
  frame=single,
  breaklines=true,
  columns=fullflexible
}

% ============================================================
% Authors and affiliations
% ============================================================
\usepackage{authblk}
\renewcommand\Affilfont{\small}
\renewcommand\Authfont{\normalsize}
\setlength{\affilsep}{0.5em}

% ============================================================
% Line numbers, headers and links
% ============================================================
\usepackage[mathlines,left]{lineno}
\usepackage{fancyhdr}
\usepackage{hyperref}
\usepackage[nameinlink,capitalise]{cleveref}

\hypersetup{
  hidelinks,
  pdftitle={Technical Note},
  pdfauthor={Author Name}
}

\pagestyle{plain} % page number centred at the bottom, as in the reference note

% ============================================================
% Bibliography
% ============================================================
\usepackage[numbers,sort&compress]{natbib}

% ============================================================
% Useful custom commands
% ============================================================
\newcommand{\nuMu}{\ensuremath{\nu_{\mu}}}
\newcommand{\anuMu}{\ensuremath{\bar{\nu}_{\mu}}}
\newcommand{\nuE}{\ensuremath{\nu_{e}}}
\newcommand{\anuE}{\ensuremath{\bar{\nu}_{e}}}
\newcommand{\GeV}{\ensuremath{\mathrm{GeV}}}
\newcommand{\MeV}{\ensuremath{\mathrm{MeV}}}
\newcommand{\todo}[1]{\textcolor{red}{[TODO: #1]}}

% Slightly compact section spacing, close to a technical-note layout
\usepackage{titlesec}
\titlespacing*{\section}{0pt}{2.2ex plus 0.8ex minus 0.2ex}{1.0ex}
\titlespacing*{\subsection}{0pt}{1.8ex plus 0.6ex minus 0.2ex}{0.7ex}

% ============================================================
% Document information
% ============================================================
\title{Technical Note: A new PID method}

\author[1]{Nicola Sommaggio}
%\author[1]{Second Author}
%\author[2]{Third Author}

\affil[1]{University of Pisa}
%\affil[2]{Second Institution or Collaboration}

\date{\today}

\begin{document}

% Continuous line numbering throughout the document
\linenumbers

\maketitle

\tableofcontents

\begin{abstract}
Abstract ...
\end{abstract}

\section{Introduction}

\section{Training set}
\subsection{Cuts}

\begin{table}
  \centering
  \begin{tabular}{c|c}
    variable & cut \\ \hline
    slice truth matching efficiency & > 0.5 \\ \hline
    $|\vec{V}^{true}-\vec{V}^{reco}|$ & < 100 cm \\ \hline
    is clear cosmic from Pandora & False \\ \hline 
  \end{tabular}
  \caption{Slice level cuts}
  \label{tab:slice_cuts_training_set}
\end{table}


\begin{table}
  \centering
  \begin{tabular}{c|c}
    variable & cut \\ \hline
    \multicolumn{2}{>{\columncolor{gray!30}}l}{\textbf{Candidate Muon}} \\ \hline
    track $S^{reco}_x$ and length & non Nan \\ \hline
    trkScore & $\geq$ 0.5 \\ \hline
    $|\vec{V}^{reco}-\vec{S}^{reco}|$ & < 10 cm \\ \hline
    track length & > 50 cm \\ \hline
    pdg code & $\pm 13$  \\ \hline
    $\vec{S}^{reco}$ contained & True \\ \hline
    $\vec{E}^{reco}$ contained & True \\ \hline
    $E^{reco}_x \cdot V^{reco}_{x}$ & > 0 \\ \hline
    from primary vertex & True \\ \hline
    longest track in the slice & True \\ \hline 
    has valid hits\tablefootnote{To have valid hits in a specific plane the track needs to have at leat 1 hit satisfying $1 \, \text{MeV/cm} < \text{dE/dx} < 30 \, \text{MeV/cm}$ and $1 \, \text{cm} < \text{RR} < 25 \, \text{cm}$ } & True \\ \hline
    truth matching completeness & $ \geq 0.5$ \\ \hline 
    truth matching efficiency & $ \geq 0.5$ \\ \hline 
    daughter electron energy match & $\leq$ 10 MeV/cm \\ \hline 
    \multicolumn{2}{>{\columncolor{gray!30}}l}{\textbf{Stopping Muon}} \\ \hline
    $|\vec{E}^{reco}-\vec{E}^{true}|$ (TODO) & $\leq$ 3 cm \\ \hline
    \multicolumn{2}{>{\columncolor{gray!30}}l}{\textbf{Muon Mip}} \\ \hline
    $|\vec{E}^{reco}-\vec{E}^{true}|$ (TODO) & $>$ 3 cm
   
  \end{tabular}
  \caption{Muon cuts}
  \label{tab:muon_cuts_trainig_set}
\end{table}


\begin{table}
  \centering
  \begin{tabular}{c|c}
    variables & cut \\ \hline
  \end{tabular}
  \caption{Proton cuts}
  \label{tab:proton_cuts_trainig_set}
\end{table}

\begin{table}
  \centering
  \begin{tabular}{c|c}
    variables & cut \\ \hline
  \end{tabular}
  \caption{Pion cuts}
  \label{tab:pion_cuts_trainig_set}
\end{table}

\end{document}
